\documentclass[a4paper, 11pt]{article}
\usepackage{geometry}
\usepackage{url}
\usepackage{color}
\usepackage[utf8]{inputenc}
\usepackage[T1]{fontenc}
\usepackage{lmodern}
\usepackage{amsmath,amssymb}
\usepackage{booktabs}
\usepackage{enumitem}


\linespread{1}

\geometry{a4paper,top=3cm,left=3cm,right=2.5cm,bottom=2cm}

\usepackage{hyperref}
\hypersetup{colorlinks,linkcolor=black,citecolor=blue}


\title{\textbf{Laboratory assignment} \\[1ex] \large \textbf{Component} (1) -- Definition of the learning task}

\author{\textbf{Authors:} Ștefan-Lucian Grămadă\\ \textbf{Group:} ICA1}

\begin{document}

\maketitle

\section{Learning task: recognition of human activities from smartphone sensor data}

The software project addresses a single learning task, \emph{supervised multi-class classification}, defined on the \emph{Human Activity Recognition Using Smartphones} data set from the UCI Machine Learning Repository~\cite{anguita2013}. The conventional ML technique chosen for this task is the Support Vector Machine~\cite{cortes1995}.

\subsection{Problem definition}

Smartphones carry inertial sensors, an accelerometer and a gyroscope, that record how the body of the person carrying the phone moves. The problem is to \textbf{automatically recognise the physical activity that a person is performing}, using only the signals recorded by a smartphone worn on the waist. The activity must be identified from a short time window of sensor readings of 2.56 seconds as one of six daily activities: \emph{walking}, \emph{walking upstairs}, \emph{walking downstairs}, \emph{sitting}, \emph{standing} and \emph{laying}.

Recognising these activities is useful in ambient assisted living and remote health monitoring, as well as in context-aware mobile applications~\cite{anguita2012}. Examples include supervising elderly or disabled people, measuring the level of physical activity and detecting long periods of inactivity.

\subsection{Problem specification}

\subsubsection{Input data}

The input of the learning system is the UCI HAR data set~\cite{anguita2013}, a labelled tabular data set with \textbf{10\,299 instances}, each characterised by \textbf{561 predefined numerical features}. The data is provided already split by volunteer into a training set and a test set.

For the learner, the input therefore consists of a set of $n$ labelled examples
\[
  D = \{(x_i, y_i)\}_{i=1}^{n}, \qquad x_i \in \mathbb{R}^{561}, \quad y_i \in C = \{1, 2, \dots, 6\},
\]
together with the volunteer identifier $s_i \in \{1,\dots,30\}$ of every example, which is used only to split the data, never as a feature.

\subsubsection{Description of the input data}

The data were collected by Anguita et al.~\cite{anguita2013} from \textbf{30 volunteers} aged 19--48 years. Each volunteer performed the six activities while wearing a Samsung Galaxy S~II smartphone on the waist. The embedded accelerometer and gyroscope recorded 3-axial linear acceleration and 3-axial angular velocity at a constant rate of 50\,Hz, and the experiments were video-recorded so that the activities could be labelled manually.

The signals were pre-processed with noise filters and split into fixed-width sliding windows of 2.56\,s with 50\% overlap, so that each window contains 128 readings. The acceleration was separated, with a low-pass filter, into a body component and a gravity component. From each window, a vector of 561 features was computed in the time domain and in the frequency domain, the latter obtained with the Fast Fourier Transform~\cite{anguita2013}. The features are derived from the following signals: body and gravity acceleration, body angular velocity, their jerk signals, which are their time derivatives, and the Euclidean magnitudes of all of them. For every signal, statistics such as the mean, standard deviation, median absolute deviation, maximum, minimum, signal magnitude area, energy, interquartile range, entropy, autoregression coefficients, correlation, skewness, kurtosis, frequency-band energies and angles between vectors were computed. A feature is named after the signal and the statistic, e.g., \texttt{tBodyAcc-mean()-X} or \texttt{fBodyGyro-energy()-Z}.

Table~\ref{tab:data} summarises the data set and Table~\ref{tab:classes} gives the class distribution.

\begin{table}[h]
\centering
\caption{Main characteristics of the UCI HAR data set.}
\label{tab:data}
\begin{tabular}{ll}
\toprule
Characteristic & Value \\
\midrule
Number of instances & 10\,299 \\
Training instances & 7\,352 \\
Test instances & 2\,947 \\
Number of features & 561 \\
Feature type & numerical, real-valued \\
Feature range & normalised to $[-1, 1]$ \\
Number of classes & 6 \\
Volunteers & 30 \\
Volunteers in the training set & 21 \\
Volunteers in the test set & 9 \\
Missing values & none \\
Instance & one 2.56\,s window \\
Readings per window & 128, sampled at 50\,Hz \\
\bottomrule
\end{tabular}
\end{table}

\begin{table}[h]
\centering
\caption{Class distribution of the UCI HAR data set.}
\label{tab:classes}
\begin{tabular}{clrrr}
\toprule
Label & Activity & Training & Test & Total \\
\midrule
1 & WALKING & 1\,226 & 496 & 1\,722 \\
2 & WALKING\_UPSTAIRS & 1\,073 & 471 & 1\,544 \\
3 & WALKING\_DOWNSTAIRS & 986 & 420 & 1\,406 \\
4 & SITTING & 1\,286 & 491 & 1\,777 \\
5 & STANDING & 1\,374 & 532 & 1\,906 \\
6 & LAYING & 1\,407 & 537 & 1\,944 \\
\midrule
 & Total & 7\,352 & 2\,947 & 10\,299 \\
\bottomrule
\end{tabular}
\end{table}

\subsubsection{Preconditions}

The following conditions must hold for the input data:
\begin{itemize}[itemsep=2pt]
  \item every instance is a vector $x \in \mathbb{R}^{561}$ with all 561 features present;
  \item every feature value lies in the interval $[-1, 1]$;
\end{itemize}

\subsubsection{Output data}

The output of the learning process is a \textbf{trained classifier} $h : \mathbb{R}^{561} \rightarrow C$, a model that approximates the unknown target function $f$ which maps every feature vector to the activity actually performed.

When applied to a new instance $x$, the classifier output the predicted activity label $\hat{y} = h(x) \in \{1,\dots,6\}$.

\subsubsection{Postconditions}

\begin{itemize}[itemsep=2pt]
  \item every input instance receives exactly one predicted label $\hat{y} \in \{1,\dots,6\}$, the class with the highest score, i.e. $\hat{y} = \arg\max_{c \in C} \mathrm{score}_c(x)$;
\end{itemize}

\subsection{Specification of the learning task (T, P, E)}
\label{sec:tpe}

\paragraph{Task (T).} Classify a 2.56-second window of smartphone sensor data, described by its 561-feature vector $x \in \mathbb{R}^{561}$, into one of the six activities $C = \{$WALKING, WALKING\_UPSTAIRS, WALKING\_DOWNSTAIRS, SITTING, STANDING, LAYING$\}$.

\paragraph{Performance measure (P).} The ability of the classifier to label correctly windows recorded from volunteers it has not seen during training. It is measured on the test data with:
\begin{itemize}[itemsep=2pt]
  \item \textbf{accuracy} -- the fraction of correctly classified instances, $acc = \frac{1}{n}\sum_{i=1}^{n} [h(x_i) = y_i]$; this is the main measure;
  \item \textbf{precision}, \textbf{recall}, also called sensitivity, and \textbf{F-measure} -- computed from the $6 \times 6$ confusion matrix for each class against all the other classes, and aggregated over the classes as a macro average and as a weighted average;
  \item \textbf{specificity}, \textbf{AUC}, the area under the ROC curve, and \textbf{AUPRC}, the area under the precision--recall curve, computed for each class with a one-vs-rest approach from the class scores and averaged over the classes.
\end{itemize}

\paragraph{Training experience (E).} A set of windows labelled by a teacher: the pairs $(x_i, y_i)$ from the training data, where $x_i$ is the 561-feature vector of a window and $y_i$ is the activity performed during that window, labelled manually from the video recordings of the experiments~\cite{anguita2013}. The experience is:
\begin{itemize}[itemsep=2pt]
  \item \emph{direct}: each training example contains the correct output of the target function for that input;
  \item \emph{with a teacher}: the labels are provided by a supervisor, as in supervised learning, and the learner does not choose its examples;
  \item \emph{representative} for the performance goal: it covers all six activities in roughly equal proportions and contains several people of different ages, so that the learned classifier can generalise to new people. In each cross-validation fold, the experience is formed by the windows of the training volunteers. In the official split, these are 7\,352 instances from 21 volunteers.
\end{itemize}

\bibliographystyle{alpha}
\bibliography{biblio}

\end{document}
